\documentclass[10pt,a4paper]{article}
\usepackage[utf8]{inputenc}
\usepackage[a4paper,left=3cm,right=2cm]{geometry}
\usepackage{amsmath}
\usepackage[thmmarks, amsmath, thref]{ntheorem}
\usepackage{amsfonts}
\usepackage{amssymb}
\usepackage{fancyhdr}
\pagestyle{fancy}
\usepackage{anysize}
\usepackage{graphicx}
\usepackage{float}
\usepackage{hyperref}
\usepackage{multirow}
\usepackage{enumitem}
\newtheorem{theorem}{Theorem}
\theoremsymbol{\ensuremath{\blacksquare}}
\newtheorem*{solution}{Solución:}

\usepackage{xcolor}

% Margenes y formalidades

\textwidth 6.25in % Margen Derecho y ancho del texto
\textheight 8.5in
\oddsidemargin 0in % Margen Izquierdo
\headheight 0in % Largo del texto
\leftmargin 1in
\parindent 0em % Salto entre la indentación (Sangría)https://www.overleaf.com/4519127669sdpptkhvsbfn
\topmargin -0.5in % Margen Superior
\parskip 0.5ex % Salto entre parrafos
\baselineskip 1pt

%\lhead{\includegraphics[scale=0.2]{logo.png}} % texto izquierda de la cabecera
%\chead{}                                      % texto centro de la cabecera
%\rhead{\includegraphics[scale=0.2]{FIG/logo-mba-2016.png}} % número de página a la derecha

\renewcommand{\headrulewidth}{0.4pt} % grosor de la línea de la cabecera
\renewcommand{\footrulewidth}{0.4pt} % grosor de la línea del pie

\begin{document}



\pagebreak
\vspace*{4mm} 

\begin{center}
\begin{huge}
\textbf{\\Ejercicios Econometría 1}
\end{Large}\\
%{\large 2023-2}\\
\end{center}


\section{Muestreo}

\textbf{Ejercicio 1:} Se tiene una población de 1000 individuos con una media de ingreso mensual de $1500$ y una desviación estándar de $300$. Si se toma una muestra aleatoria simple de 100 individuos, ¿cuál es la media esperada y la desviación estándar de la media muestral?

\textbf{Solución:}

La media esperada de la muestra es igual a la media poblacional:
$$E(\Bar{X}) = \mu = 1500$$


La desviación estándar de la media muestral (error estándar) se calcula como:

$$\sigma_{\Bar{X}} = \frac{\sigma}{\sqrt{n}} = \frac{300}{\sqrt{100}} = \frac{300}{10} = 30$$

\section{Estadística Descriptiva}

\textbf{Ejercicio 2:} Se tiene el siguiente conjunto de datos sobre el número de ventas realizadas por 5 vendedores en un mes:
$$X = \{10,15,12,18,20\}$$
Calcule la media, mediana y varianza muestral.

\textbf{Solución:}

La media es:
$$\Bar{X} = \frac{10+15+12+18+20}{5} = \frac{75}{5} = 15$$

Ordenando los datos: {10, 12, 15, 18, 20}, la mediana es el valor central, es decir:

$$mediana(X) = 15$$

La varianza muestral se calcula como:

$$s^2 = \frac{\sum (X_i - \Bar{X})^2}{n-1}$$

$$s^2 = \frac{(10-15)^2 + (15-15)^2 + (12 - 15)^2 + (18-15)^2 + (20-15)^2}{4}$$

$$s^2 = \frac{(-5)^2+(0)^2 + (-3)^2 +(3)^2 + (5)^2}{4} = \frac{25+0+9+9+25}{4} = \frac{68}{4} = 17$$



\section{Regresión Lineal Simple}

\textbf{Ejercicio 3:} Se tiene la siguiente muestra de datos:
\begin{center}
\begin{tabular}{|c|c|}
\hline
X (Publicidad) & Y (Ventas) \\
\hline
2 & 4 \\
\hline
4 & 7 \\
\hline
6 & 10 \\
\hline
8 & 13 \\
\hline
10 & 16 \\
\hline
\end{tabular}
\end{center}

Ajuste un modelo de regresión lineal simple de la forma:
$$Y = \beta_0 + \beta_1 X + \epsilon $$

usando el método de mínimos cuadrados ordinarios.

\textbf{Solución:}

Los estimadores de MCO se calculan como:
$$\hat{\beta}_1 = \frac{\sum (X_i - \Bar{X})(Y_i - \Bar{Y})}{\sum (X_i - \Bar{X})^2} = \frac{Cov(X,Y)}{Var(X)} $$
$$\hat{\beta}_0  = \Bar{Y} - \hat{\beta}_1 \Bar{X}$$

Calculamos las medias:
$$\Bar{X} = \frac{2+4+6+8+10}{5}=6$$ 
$$\Bar{Y} = \frac{4+7+10+13+16}{5} = 10$$

Calculamos $\hat{\beta}_1$ :
$$\hat{\beta}_1 = \frac{(2-6)(4-10)+(4-6)(7-10)+(6-6)(10-10)+(8-6)(13-10) + (10-6)(16-10)}{(2-6)^2+(4-6)^2+(6-6)^2+(8-6)^2+(10-6)^2}$$
$$\hat{\beta}_1 = \frac{(-4)(-6)+(-2)(-3)+(0)(0)+(2)(3)+(4)(6)}{(-4)^2+(-2)^2+(0)^2+(2)^2+(4)^2}$$
$$\hat{\beta}_1 = \frac{24+6+0+6+24}{16+4+0+4+16}= \frac{60}{40} = 1.5$$

Calculamos $\hat{\beta}_0$:
$$\hat{\beta}_0 = \Bar{Y} - \hat{\beta}_1 \Bar{X} = 10-(1.5)6 = 10 -9 = 1$$

Por lo tanto, la ecuación estimada es:

$$\hat{Y} = 1+1.5X$$

\section{Regresión Lineal Simple: Educación y Salario}

\textbf{Ejercicio:} Se tiene la siguiente muestra de datos sobre los años de educación y el salario mensual (en dólares) de 10 individuos:

\begin{center}
\begin{tabular}{|c|c|}
\hline
Años de Educación ($X$) & Salario Mensual ($Y$) \\
\hline
12  & 600  \\
\hline
13  & 750  \\
\hline
14  & 900  \\
\hline
15  & 1100 \\
\hline
16  & 1300 \\
\hline
18  & 1600 \\
\hline
20  & 1850 \\
\hline
22  & 2100 \\
\hline
24  & 2300 \\
\hline
25  & 2500 \\
\hline
\end{tabular}
\end{center}

Ajuste un modelo de regresión lineal simple de la forma:
$$Y = \beta_0 + \beta_1 X + \epsilon $$


usando el método de mínimos cuadrados ordinarios, con el cálculo explícito de la varianza y la covarianza mediante tablas.

\textbf{Solución:}

Calculamos las medias:
$$\Bar{X} = \frac{12+13+14+15+16+18+20+22+24+25}{10} = 17.9$$
$$\Bar{Y} = \frac{600+750+900+1100+1300+1600+1850+2100+2300+2500}{10} = 1500$$

Construimos la siguiente tabla:

\begin{center}
\begin{tabular}{|c|c|c|c|c|c|}
\hline
& & & & &\\
$X$ & $Y$ & $(X - \Bar{X})$ & $(Y - \Bar{Y})$ & $(X - \Bar{X})(Y - \Bar{Y})$ & $(X - \bar{X})^2$ \\
\hline
12  & 600  & -5.9  & -900  & 5310  & 34.81 \\
\hline
13  & 750  & -4.9  & -750  & 3675  & 24.01 \\
\hline
14  & 900  & -3.9  & -600  & 2340  & 15.21 \\
\hline
15  & 1100 & -2.9  & -400  & 1160  & 8.41  \\
\hline
16  & 1300 & -1.9  & -200  & 380   & 3.61  \\
\hline
18  & 1600 & 0.1   & 100   & 10    & 0.01  \\
\hline
20  & 1850 & 2.1   & 350   & 735   & 4.41  \\
\hline
22  & 2100 & 4.1   & 600   & 2460  & 16.81 \\
\hline
24  & 2300 & 6.1   & 800   & 4880  & 37.21 \\
\hline
25  & 2500 & 7.1   & 1000  & 7100  & 50.41 \\
\hline
& & & total & 28050 & 194.9 \\
\hline
\end{tabular}
\end{center}

Ahora calculamos la covarianza y la varianza:
$$Cov(X,Y) = \frac{\sum (X_i - \Bar{X})(Y_i - \Bar{Y})}{n-1}= \frac{28050}{9} = 3116.67$$
$$Var(X) = \frac{\sum (X_i-\Bar{X})^2}{n-1} = \frac{194.9}{9} = 21.65$$

Calculamos $\hat{\beta}_0$ y $\hat{\beta}_1$:
$$\hat{\beta}_1 = \frac{Cov(X,Y)}{Var(X)} = \frac{3116.67}{21.65} = 143.9$$
$$\hat{\beta}_0 = \Bar{Y} - \hat{\beta}_1 \Bar{X} = 1500-(143.9)(17.9) = -1071.85$$

Por lo tanto, la ecuación estimada es:
$$\hat{Y} = -1071.85 + 143.9 X$$

\section{Regresión Lineal Simple: Descomposición de Varianza y $R^2$}

\textbf{Ejercicio:} Se tiene la siguiente muestra de datos sobre la experiencia laboral (en años) y el ingreso mensual (en dólares) de 7 individuos:

\begin{center}
\begin{tabular}{|c|c|}
\hline
Experiencia Laboral ($X$) & Ingreso Mensual ($Y$) \\
\hline
2   & 1400  \\
\hline
4   & 1800  \\
\hline
6   & 2200  \\
\hline
8   & 2600  \\
\hline
10  & 3100  \\
\hline
12  & 3600  \\
\hline
14  & 4200  \\
\hline
\end{tabular}
\end{center}

Ajuste un modelo de regresión lineal simple de la forma:

$$Y = \beta_0 + \beta_1 X + \epsilon $$

usando el método de mínimos cuadrados ordinarios y calcule la descomposición de la varianza: SST, SSE y SSR, junto con el coeficiente de determinación $R^2$.

\textbf{Solución:}

Calculamos las medias:
$$\Bar{X} = \frac{2+4+6+8+10+12+14}{7}=8$$
$$\Bar{Y} = \frac{1400+1800+2200+2600+3100+3600+4200}{7} = 2700$$

Calculamos la varianza y covarianza utilizando la siguiente tabla:

\begin{center}
\begin{tabular}{|c|c|c|c|c|c|}
\hline
& & & & & \\
$X$ & $Y$ & $(X - \bar{X})$ & $(Y - \bar{Y})$ & $(X - \bar{X})^2$ & $(X - \bar{X})(Y - \bar{Y})$ \\
\hline
2   & 1400  & -6  & -1300  & 36   & 7800  \\
\hline
4   & 1800  & -4  & -900   & 16   & 3600  \\
\hline
6   & 2200  & -2  & -500   & 4    & 1000  \\
\hline
8   & 2600  & 0   & -100   & 0    & 0     \\
\hline
10  & 3100  & 2   & 400    & 4    & 800   \\
\hline
12  & 3600  & 4   & 900    & 16   & 3600  \\
\hline
14  & 4200  & 6   & 1500   & 36   & 9000  \\
\hline
 & & & Total & 112 & 25800 \\
 \hline
\end{tabular}
\end{center}

Ahora calculamos la varianza y la covarianza:
$$Var(X) = \frac{\sum (X_i -\Bar{X})^2}{n-1} = \frac{112}{6} = 18.67$$
$$Cov(X,Y) = \frac{\sum (X_i  - \Bar{X})(Y_i - \Bar{Y})}{n-1} = \frac{25800}{6} = 4300$$

Calculamos los coeficientes de la regresión:
$$\hat{\beta}_1 = \frac{Cov(X,Y)}{Var(X)} = \frac{4300}{18.67} = 230,32$$
$$\hat{\beta}_0 = \Bar{Y}-\hat{\beta}_1 \Bar{X} = 2700-(230.32)(8)=857,44$$

Por lo tanto, la ecuación estimada es:
$$\hat{Y} = 857.44 + 230.32 X$$

Ahora detallamos el cálculo de SST, SSE y SSR en la siguiente tabla:

\begin{center}
\begin{tabular}{|c|c|c|c|c|c|}
\hline
$X$ & $Y$ & $\hat{Y}$ & $(Y - \bar{Y})^2$ & $(\hat{Y} - \bar{Y})^2$ & $(Y - \hat{Y})^2$ \\
\hline
2   & 1400  & 1318,08  & 1690000  & 1909702,89  & 6710,89  \\
\hline
4   & 1800  & 1778.72  & 810000   & 848756,84   & 452,84  \\
\hline
6   & 2200  & 2239,36  & 250000   & 212189,21   & 1549,21  \\
\hline
8   & 2600  & 2700 & 10000    & 0        & 10000  \\
\hline
10  & 3100  & 3160,64  & 160000   & 212189,21   & 3677,21  \\
\hline
12  & 3600  & 3621,28  & 810000   & 452,84   & 452,84  \\
\hline
14  & 4200  & 4081,92  & 2250000  & 1909702,89  & 13942,89  \\
\hline
\end{tabular}
\end{center}

Finalmente, sumamos las columnas para obtener:

$$SST = \sum (Y_i - \Bar{Y})^2=8360000$$

$$SSE = \sum (\hat{Y}_i - \Bar{Y})^2= 7787900$$

$$SSR = \sum (Y_i - \hat{Y}_i)^2 = 572100$$

Verificamos la relación fundamental:

$$SST = SST+SSE \implies 8360000=7787900+572100$$

El coeficiente de determinación $R^2$ se calcula como:

$$R^2 = \frac{SSE}{SST} = \frac{7787900}{8360000} = 0.93$$


Esto indica que el 93\% de la variabilidad del ingreso mensual es explicada por la experiencia laboral.

\end{document}